\section*{Chapter \ref{ch:b2:himalaya-andes} --- Collision and Andean Belts Compared}

\begin{solution}{exo:b2:himalaya-andes:1}
(a) $r = h\rho_c/(\rho_m - \rho_c) = 3.0\times2700/600 = \qty{13.5}{km}$. (b) $T = 35 + 3.0 +
13.5 = \qty{51.5}{km}$, or $35 + 5.5\times3.0$. (c) $h = \Delta T(\rho_m - \rho_c)/\rho_m =
10\times600/3300 = \qty{1.8}{km}$.
\end{solution}

\begin{solution}{exo:b2:himalaya-andes:2}
(a) Indo-Gangetic foredeep; Main Frontal Thrust; sub-Himalaya (Siwaliks); Main Boundary
Thrust; Lesser Himalaya; Main Central Thrust; Greater Himalayan sequence; South Tibetan
Detachment System; Tethyan Himalayan sequence; then the Indus--Tsangpo suture and the Asian
crust. (b) The South Tibetan Detachment System. (c) The grade rises northwards: unmetamorphosed
\omterm{def:b2:alpine-orogen:foreland}{molasse}, medium-grade Lesser Himalaya, high-grade and partly molten Greater Himalayan
sequence; above the STDS it drops abruptly to unmetamorphosed or weakly metamorphosed
Tethyan sediments, since the normal fault has removed the rocks in between.
\end{solution}

\begin{solution}{exo:b2:himalaya-andes:3}
(a) $196.5\ \text{m}/9200\ \text{yr} = \qty{0.02136}{m}$ a year, about \qty{21.4}{mm} a year.
(b) It is \qty{20}{\percent} more than the \qty{17.8}{mm} a year of GPS, but of the same size:
one rate is averaged over nine thousand years, the other over twenty.
(c) Nearly all the shortening across the range is taken up at the front: the MBT, the MCT and
the other faults north of the MFT are, at this longitude, practically inactive; they all root
in the same MHT, which comes up at the MFT.
\end{solution}

\begin{solution}{exo:b2:himalaya-andes:4}
(a) Collision: the suture marks the vanished ocean. (b) Andean-type: an arc needs a slab
going down. (c) Collision, typically: crustal melts of a thickened, heated crust (they also
occur in thickened Andean crust, but the arc \omterm{def:b2:magma-emplacement:magma}{magmas} dominate there). (d) Andean-type.
(e) Both: the Siwaliks and the Subandean foreland. (f) Both: the MHT in the Himalaya, the
subduction interface in the Andes.
\end{solution}

\begin{solution}{exo:b2:himalaya-andes:5}
(a) $17.8/45 = \qty{40}{\percent}$ in the east, $20.5/45 = \qty{46}{\percent}$ in the west:
about half. (b) Into the thickening and the faults of Tibet and the rest of Asia, and into
the eastward extrusion of blocks along the great strike-slip faults. (c) $45 - 17.8 =
\qty{27.2}{mm}$ a year, or \qty{272}{km} in \qty{10}{Myr}; $45 - 20.5 = \qty{24.5}{mm}$ a
year, or \qty{245}{km}.
\end{solution}

\begin{solution}{exo:b2:himalaya-andes:6}
(a) $r = 4.5\times3.7 = \qty{16.65}{km}$; $\Delta F = 2700\times9.8\times3700\times(35\,000 +
1850 + 8325) = \qty{4.42e12}{N.m^{-1}}$. (b) Tibet: \qty{6.12e12}{N.m^{-1}}; the Altiplano
has \qty{72}{\percent} of it. (c) With $r = 4.5h$, $\Delta F = \rho_c g\,(Hh + 2.75h^2)$: for
$h \ll H$ the term $Hh$ dominates and $\Delta F$ is proportional to $h$; for a high plateau
the term in $h^2$ grows faster. Each further kilometre of height costs more and more push at
the edges, until the forces available can no longer raise the plateau; it then spreads
instead of rising.
\end{solution}

\begin{solution}{exo:b2:himalaya-andes:7}
(a) $2700\times9.8\times8848.86/2 = \qty{1.17e8}{Pa} = \qty{117}{MPa}$. (b) At
\qty{100}{MPa}: $0.85\times100 = \qty{85}{MPa}$; at \qty{200}{MPa}: $0.85\times200 = 50 +
0.6\times200 = \qty{170}{MPa}$, the two branches agree. (c) $h = 2\sigma_d/(\rho g) =
2\times1.70\times10^8/(2700\times9.8) = \qty{12.8}{km}$: an order of magnitude, of the order
of the highest summits, not far above them.
\end{solution}

\begin{solution}{exo:b2:himalaya-andes:8}
(a) $11\times111.19\times\cos\ang{30} = \qty{1059}{km}$. (b) $9\times10^{-3}/1.059\times10^{6}
= \num{8.5e-9}$ a year (between \num{6.6e-9} and \num{10.4e-9} for 7 and \qty{11}{mm} a year),
that is $8.5\times10^{-9}/3.156\times10^{7} = \qty{2.7e-16}{s^{-1}}$. (c) The two ranges
overlap: $(12.4 \pm 4.7) \times 10^{-9}$ spans 7.7 to \num{17.1e-9}. The southern plateau north of
Nepal stretches as fast as, or a little faster than, the belt as a whole.
\end{solution}

\begin{solution}{exo:b2:himalaya-andes:9}
(a) $0.9/6.6 = \qty{14}{\percent}$. (b) $M_0 = 10^{1.5\times8.5 + 9.1} = 10^{21.85} =
\qty{7.1e21}{N.m}$. (c) $7.1\times10^{21}/6.6\times10^{19} = 107$: one such earthquake about
every 110 years beneath Nepal, far more than history records; either much larger earthquakes
are due, or part of the deficit is released otherwise.
\end{solution}

\begin{solution}{exo:b2:himalaya-andes:10}
(a) Take the compensation depth at the base of the Altiplano crust, $61.5 - 3.7 =
\qty{57.8}{km}$ below sea level. The Puna crust reaches $50.5 - 4.4 = \qty{46.1}{km}$, so the
Puna column holds \qty{11.7}{km} of mantle above that depth: $2800\times61.5 =
\rho_P\times50.5 + 3300\times11.7$. (b) $\rho_P = (172\,200 - 38\,610)/50.5 =
\qty{2645}{kg.m^{-3}}$. (c) The difference, \qty{155}{kg.m^{-3}}, lies at the lower end of
the 150 to \qty{200}{kg.m^{-3}} quoted: crust alone would need an unusually light Puna crust.
A lighter column beneath it does it more easily: a mantle lithosphere removed by delamination,
replaced by hot, light asthenosphere.
\end{solution}

\begin{solution}{exo:b2:himalaya-andes:11}
(a) $g = 105/\tan\ang{30} = \qty{181.9}{km}$. (b) $g = 90/\tan\ang{30} + 15/\tan\ang{4} =
155.9 + 214.5 = \qty{370.4}{km}$: the arc moves \qty{188.5}{km} inland. (c) The slab never
reaches \qty{105}{km} beneath the plate for those \qty{300}{km}, and no hot mantle wedge lies
above it: no \omterm{def:b2:magma-emplacement:magma}{magma} is made and the arc dies until the slab sinks again, further east.
\end{solution}

\begin{solution}{exo:b2:himalaya-andes:12}
(a) \qty{21}{mm} a year is \qty{21}{km} per million years: $21\times37.5 =
\qty{787.5}{km^3}$ per kilometre of range per million years. (b) $300/787.5 =
\qty{38}{\percent}$. (c) In the thickening of the range and of southern Tibet, in crust
carried down with India beneath Tibet, and in the material that the MHT passes northwards.
\end{solution}

\begin{solution}{pb:b2:himalaya-andes:1}
\textbf{1.} Seven samples from \ang{28.0} to \ang{37.0}~N: $34\,550/7 = \qty{4936}{m}$, $h_T
= \qty{4.9}{km}$.
\textbf{2.} Five samples from \ang{-69} to \ang{-65}: $18\,650/5 = \qty{3730}{m}$, $h_A =
\qty{3.7}{km}$, the value quoted for the Altiplano; the Puna, further south, is higher.
\textbf{3.} Tibet: at least $9\times111.19 = \qty{1001}{km}$; the Altiplano: at least $4\times
104.5 = \qty{418}{km}$.
\textbf{4.} $(5000 - 70)/(1.5\times111\,190) = 0.0296$: \qty{3.0}{\percent}.
\textbf{5.} The Andean profile starts on the ocean floor and in the trench, kilometres below
sea level, west of the plateau; the Tibetan one starts on a continental plain near sea level,
and ends on the high Tarim basin.
\textbf{6.} $r = 4.5\times4.9 = \qty{22.05}{km}$; $T_T = 35 + 4.9 + 22.05 = \qty{62.0}{km}$.
\textbf{7.} $r = 4.5\times3.7 = \qty{16.65}{km}$; $T_A = 35 + 3.7 + 16.65 = \qty{55.4}{km}$.
\textbf{8.} Tibet: \qty{62.0}{km} against 70 to \qty{75}{km}, \qty{10.5}{km} short of the mid
value \qty{72.5}{km} (\qty{14}{\percent}); the Altiplano: \qty{55.4}{km} against 59 to
\qty{64}{km}, \qty{6.1}{km} short of \qty{61.5}{km} (\qty{10}{\percent}). Tibet is further off.
\textbf{9.} $72.5 - 35 = 4.9\times3300/(3300 - \rho_c)$: $3300 - \rho_c = 16\,170/37.5 =
431.2$, $\rho_c = \qty{2869}{kg.m^{-3}}$.
\textbf{10.} $26.5 = 3.7\times3300/(3300 - \rho_c)$: $3300 - \rho_c = 460.8$, $\rho_c =
\qty{2839}{kg.m^{-3}}$.
\textbf{11.} Part of the load is carried by the bending of the plate, not by a local root; the
normal crust may not have been \qty{35}{km} thick; the mantle beneath may be lighter or denser
than assumed; the seismic Moho may not be the base of the crust in the density sense.
\textbf{12.} The base of Tibet's crust lies at $72.5 - 4.9 = \qty{67.6}{km}$, India's at
\qty{37.5}{km}: $t = \qty{30.1}{km}$. $\Delta g = 2\pi\times6.674\times10^{-11}\times370\times
30\,100 = \qty{4.67e-3}{m.s^{-2}}$, that is \qty{-467}{mGal} relative to India, close to the
observed \qty{460}{mGal}.
\textbf{13.} $W_0 = 440\times60/35 = \qty{754.3}{km}$; $S = \qty{314.3}{km}$.
\textbf{14.} $(314.3 - 250)/314.3 = \qty{20}{\percent}$, the shortfall found in the Central
Andes.
\textbf{15.} $0.55\times314.3 = \qty{173}{km}$, less than the \qty{250}{km} of the balanced
sections, which would then be more than enough.
\textbf{16.} The cover shortened above the \omterm{def:b2:alpine-orogen:decollement}{décollement} must be matched by shortening of the
basement beneath it, further west, since both must lose the same length across the belt; the basement
shortening thickens the crust under the Eastern Cordillera and the plateau.
\textbf{17.} $314.3/30 = \qty{10.5}{mm}$ a year, $10.5/72 = \qty{15}{\percent}$ of the
convergence.
\textbf{18.} It is consumed by subduction: the Nazca plate sinks beneath the continent.
\textbf{19.} $r = \qty{22.05}{km}$; $\Delta F = 2700\times9.8\times4900\times(35\,000 + 2450
+ 11\,025) = \qty{6.28e12}{N.m^{-1}}$.
\textbf{20.} $2700\times9.8\times3700\times45\,175 = \qty{4.42e12}{N.m^{-1}}$; the ratio is
0.70.
\textbf{21.} $6.28\times10^{12}/10^{5} = \qty{6.3e7}{Pa}$, about \qty{63}{MPa}.
\textbf{22.} $9\times10^{-3}/1.06\times10^{6} = \num{8.5e-9}$ a year.
\textbf{23.} The push of India acts north--south and is balanced, in the high plateau, by the
plateau's own weight: north--south the plateau can no longer thicken, so the vertical stress,
the weight, is the greatest; the least stress is east--west, towards the free edges of Asia,
and the plateau extends that way, on north--south rifts.
\textbf{24.} \textbf{\qty[text-series-to-math=true]{62}{km}}, about \qty{86}{\percent} of the
\qty{72.5}{km} measured.
\end{solution}
